\documentclass[11pt]{article}

\usepackage[margin=1in]{geometry}
\usepackage{parskip}
\usepackage{hyperref}
\usepackage{xcolor}
\usepackage{listings}
\usepackage{booktabs}
\usepackage{longtable}
\usepackage{enumitem}
\usepackage{amsmath}

\hypersetup{colorlinks=true, linkcolor=blue, urlcolor=blue, citecolor=blue}

\definecolor{codebg}{rgb}{0.96,0.96,0.96}
\lstset{
  basicstyle=\ttfamily\small,
  backgroundcolor=\color{codebg},
  breaklines=true,
  columns=fullflexible,
  keywordstyle=\color{blue!70!black}\bfseries,
  commentstyle=\color{green!40!black},
  stringstyle=\color{purple},
  frame=single,
  framesep=3pt,
  xleftmargin=3pt,
  xrightmargin=3pt,
  tabsize=2,
}
\lstdefinelanguage{PyShort}{language=Python}

\title{How AutoNixNan Works: A Code-Level Account}
\author{ganesh@cs.utah.edu \\ \small (written with Claude Code)}
\date{\today}

\begin{document}
\maketitle

\begin{abstract}
This document explains, from the source code, how \texttt{bin/autonixnan}
(``AutoNixNan'') and the \texttt{nixnan.so} NVBit tool it drives actually
work. It is not a user guide; it is a line-level account of the
mechanism, written after reading every file involved
(\texttt{bin/autonixnan}, \texttt{src/nixnan.cu}, and the eleven
supporting headers/translation units under \texttt{src/}), and after
empirically exercising the tool against several real CUDA programs. Where
the code's behavior diverges from what its own comments or the project
README claim, that divergence is called out explicitly and is backed by
a repeatable observation, not speculation.
\end{abstract}

\tableofcontents

\section{The Big Picture}

AutoNixNan is not one program; it is two, layered:

\begin{enumerate}
\item \textbf{\texttt{nixnan.so}} (source: \texttt{src/*.cu}, \texttt{src/*.cuh},
  \texttt{src/*.h}) --- an NVBit tool. NVBit (NVIDIA Binary Instrumentation
  Tool) is preloaded into a CUDA process via \texttt{LD\_PRELOAD} and gets
  callbacks whenever the process launches a kernel, before that kernel's
  SASS (compiled machine code) ever executes on the GPU. \texttt{nixnan.so}
  uses those callbacks to rewrite each kernel's SASS, inserting calls to
  its own device-side functions after (and sometimes before) every
  floating-point instruction, so it can inspect the actual bits flowing
  through the GPU's registers at run time.
\item \textbf{\texttt{bin/autonixnan}} (a 420-line Python 3 script) --- a
  wrapper that never touches SASS or the GPU directly. It only launches
  the target program multiple times, each time with a different
  combination of environment variables that steer \texttt{nixnan.so},
  and stitches the resulting log files into one report.
\end{enumerate}

Everything \texttt{autonixnan} does, it does by shelling out to the
target program with \texttt{LD\_PRELOAD=nixnan.so} and a specific
environment. There is no IPC between the Python process and the GPU
process beyond environment variables in and a text log file out. This
matters: every quirk of \texttt{autonixnan}'s behavior traces back to
either (a) what environment variable combination it chose to set for a
given phase, or (b) what \texttt{nixnan.so} actually does --- which is
sometimes \emph{not} what its own \texttt{--help} banner or
\texttt{autonixnan}'s comments say it does (Section~\ref{sec:timeout-dead}).

\section{Part I --- \texttt{nixnan.so}}

\subsection{Initialization: \texttt{nvbit\_at\_init}}

NVBit calls \texttt{nvbit\_at\_init()} once, before the target's
\texttt{main()} runs, but after CUDA driver initialization has begun
enough that reading environment variables is possible. \texttt{nixnan.cu}
uses this hook to read every environment variable it supports via the
NVBit-provided \texttt{GET\_VAR\_INT}/\texttt{GET\_VAR\_STR} macros
(defined in \texttt{nvbit\_release/core/nvbit.h}). Each call to one of
these macros \emph{also prints a banner line} to the log stream --- this
is the ``NVBit core environment variables'' banner every instrumented run
of this session emitted. That banner is not decorative: it is a direct
enumeration of every env var the running binary actually reads, in the
order it reads them. \textbf{If a variable does not appear in the banner,
the running binary does not consult it, full stop} --- this becomes the
key diagnostic in Section~\ref{sec:timeout-dead}.

Table~\ref{tab:envvars} lists every variable read by \texttt{nixnan.cu}
itself, in banner order, plus two read by a header it includes
(\texttt{fp-histogram.cu}) during \texttt{nvbit\_at\_init}'s call to
\texttt{nixnan::fp\_histogram::init()}.

\begin{longtable}{p{3.2cm}p{2.3cm}p{8.2cm}}
\toprule
\textbf{Variable} & \textbf{Default} & \textbf{Effect} \\
\midrule
\endhead
\texttt{INSTR\_BEGIN} / \texttt{INSTR\_END} & 0 / UINT32\_MAX & Instruction-index interval to instrument (unused by autonixnan). \\
\texttt{TOOL\_VERBOSE} & 0 & Extra diagnostic printing during instrumentation setup. \\
\texttt{ENABLE\_FUN\_DETAIL} & 0 & Print the full mangled kernel name every time a kernel re-runs, not just the first time. \\
\texttt{PRINT\_ILL\_INSTR} & 0 & (declared, not read further in the code paths exercised here). \\
\texttt{SAMPLING} & 0 & Instrument only every $N$th launch of a repeated kernel. \\
\texttt{INSTRUMENT\_EXCEPTIONS} & 1 & Master switch for the NaN/Inf/Subnormal/Div-by-0 instrumentation described in Section~\ref{sec:exceptions}. \\
\texttt{INSTR\_MEM} & 0 & Also instrument 32/64/128-bit global-memory stores (\texttt{STG}) for NaN, Section~\ref{sec:memcheck}. \\
\texttt{HISTOGRAM} & 0 & Set implicitly to 1 if \texttt{BIN\_SPEC\_FILE} is non-empty; enables the exponent-histogram/bin-threshold mechanism, Section~\ref{sec:histogram}. \\
\texttt{BIN\_SPEC\_FILE} & (empty) & Path to a JSON file describing which exponent ranges, per format, should be watched; auto-created as a template and the process exits if the file does not yet exist. \\
\texttt{LOGFILE} & (empty) & Redirect all \texttt{nnout()} output (everything prefixed \texttt{\#nixnan:}) from \texttt{stderr} to this file. \\
\texttt{LINE\_INFO} & 1 & Attempt to resolve source file:line for each instrumented instruction (wrapped in a \texttt{SIGSEGV} handler --- see \texttt{recording.cu}). \\
\texttt{LOG\_KERNELS} & (empty) & Path to a log file; if set, \textbf{disables all instrumentation} and instead just logs every kernel launch and its wall-clock execution time. Section~\ref{sec:logkernels}. \\
\texttt{MAX\_ERRORS} & 0 (unlimited) & Self-terminate (\texttt{SIGINT} via \texttt{nnterminate}) after this many exception reports. \\
\texttt{TIME\_KERNELS} & 0 & Print a ``Kernel [\ldots] execution time: N microseconds'' line for every kernel launch, independent of \texttt{LOG\_KERNELS}. \\
\texttt{FUNCTION\_WHITELIST} & (empty) & Path to a newline-separated file of function names; only these are instrumented. \\
\texttt{FUNCTION\_BLACKLIST} & (empty) & Complement of the above; mutually exclusive with it (the process exits with an error if both are set). \\
\bottomrule
\caption{Environment variables consulted by the running \texttt{nixnan.so} binary, in the order \texttt{nvbit\_at\_init} reads them. Absence from this table (e.g.\ \texttt{NIXNAN\_TIMEOUT}) means the variable is \emph{not currently wired up} -- see Section~\ref{sec:timeout-dead}.}
\label{tab:envvars}
\end{longtable}

\subsection{Deciding what to instrument}

For every kernel launch, NVBit calls
\texttt{nvbit\_at\_cuda\_event(ctx, is\_exit, cbid, \ldots)}, once on entry
and once on exit, for the CUDA driver calls that launch work
(\texttt{cuLaunchKernel} and its cooperative variants). On entry
(\texttt{is\_exit == false}):

\begin{lstlisting}[language=C++]
bool enable_instr = instrument_function(ctx, p->f);
if (enable_instr) {
  int count = analyzed_kernels[short_name]++;
  if (count == 0) nnout() << "running kernel [" << short_name << "] ...";
}
nvbit_enable_instrumented(ctx, p->f, true);
\end{lstlisting}

\texttt{instrument\_function} walks the kernel plus every ``related
function'' NVBit knows it might call (\texttt{nvbit\_get\_related\_functions}
--- this includes every template instantiation the compiler linked in as a
\emph{possible} callee, which is why a single cuBLAS \texttt{nrm2\_kernel}
launch can enumerate dozens of unrelated-looking function names; see
Section~\ref{sec:logkernels}). For each such function it calls
\texttt{should\_instrument}, which is where whitelist/blacklist filtering,
sampling, and --- critically --- \texttt{LOG\_KERNELS} mode all converge:

\begin{lstlisting}[language=C++]
bool should_instrument(CUcontext ctx, CUfunction f) {
  ...
  enable_instr &= !kernel_logging_enabled;   // <-- LOG_KERNELS forces this to false
  if (sampling != 0 && analyzed_kernels.count(func_name)) {
    if (analyzed_kernels[func_name] % sampling != 0) enable_instr = false;
  }
  return enable_instr;
}
\end{lstlisting}

So \texttt{LOG\_KERNELS} does not merely add logging on top of the normal
instrumentation --- it hard-disables \texttt{should\_instrument} for
\emph{every} function, meaning \texttt{instrument\_function} inserts zero
device-side calls anywhere, and the ``running kernel [\ldots]'' print above
never fires either (since \texttt{enable\_instr} is false). The kernel
still runs, at essentially native speed, because
\texttt{nvbit\_enable\_instrumented} is called unconditionally and simply
makes NVBit run the (uninstrumented, in this case identical-to-original)
version.

For every instruction NVBit hands back for a function that \emph{is}
selected, \texttt{instruction\_info::get\_reginfo(instr)} decides whether
the instruction is floating-point-interesting at all, and if so, what its
operands look like.

\subsection{The opcode-to-operand table}

\texttt{get\_reginfo} (in \texttt{instruction\_info.cu}) matches the
instruction's opcode string against a table of regex patterns baked into
the binary as \texttt{src/instructions.h} --- a byte array
(\texttt{instructions\_json}, 5{,}433 bytes) that is really just a JSON
document, embedded at compile time and parsed once at first use with
\texttt{nlohmann::json}. Two representative entries:

\begin{lstlisting}
"HMMA.16816.F16": {
  "registers": [
    {"type": "f16", "count": 4},
    {"type": "f16", "count": 8},
    {"type": "f16", "count": 4},
    {"type": "f16", "count": 4}
  ],
  "description": "Half-precision matrix multiply accumulate into half-precision accumulator"
},
"DFMA(\\.RZ|\\.RP|\\.RM)?(\\.FTZ)?(?:\\.(RP|RM|RZ))?(?:\\.SAT)?": {
  "registers": [
    {"type": "f64", "count": 1},
    {"type": "f64", "count": 1},
    {"type": "f64", "count": 1}
  ],
  "description": ""
}
\end{lstlisting}

Each entry says, per operand position, which floating-point format the
operand is (\texttt{f16}/\texttt{bf16}/\texttt{f32}/\texttt{f64}) and how
many scalar values are packed into it (e.g.\ a tensor-core \texttt{HMMA}
operand carrying 8 packed half-precision values). For each such operand,
\texttt{get\_regs} figures out \emph{how} to actually read the value at
instrumentation-insertion time --- a register, a uniform register, or a
constant-bank (\texttt{cbank}) immediate --- and builds a list of
zero-argument closures (\texttt{reginsertion}, i.e.\
\texttt{std::function<void()>}) that, when called later, will emit the
NVBit call
(\texttt{nvbit\_add\_call\_arg\_reg\_val}/\texttt{\_ureg\_val}/\texttt{\_cbank\_val})
that actually smuggles that operand's bits into the injected device
function call. A literal double/float immediate operand
(\texttt{IMM\_DOUBLE}) is instead classified right there at instrumentation
time via host-side \texttt{fpclassify} and, if already NaN/Inf/subnormal,
produces an immediate warning line rather than a runtime check (there is
nothing to check at run time --- the value never changes).

This table is also where the 24-line \texttt{is\_f64high} check matters:
some SASS forms split a 64-bit double across two 32-bit registers and
give the \emph{high} half its own opcode suffix (\texttt{...64H}); the
table's register-start-offset logic (\texttt{reg\_start = op->u.reg.num -
(f64high ? 1 : 0)}) exists purely to correctly pair the two halves back
into one 64-bit value.

\subsection{Exception detection, end to end}
\label{sec:exceptions}

For every FP-classified instruction (and \texttt{INSTRUMENT\_EXCEPTIONS} is
1 by default), \texttt{instrument\_function} in \texttt{nixnan.cu} inserts
\emph{two} calls to the same device function, \texttt{nixnan\_check\_regs}
(defined in \texttt{inject\_funs.cu}): once \texttt{IPOINT\_AFTER} (reading
the operand \emph{after} it takes its post-instruction guard-predicate
value --- effectively the output/destination operands) and once
\texttt{IPOINT\_BEFORE} (the input operands). Each call carries, as
variadic arguments, a packed \texttt{reginfo} struct (a bit-packed 32-bit
struct --- \texttt{count:6, type:4, half\_h0:1, half\_h1:1, div0:1,
operand:2, num\_regs:4, f64high:1}, statically asserted to fit in one
\texttt{int32\_t}) followed by the actual register values.

\texttt{nixnan\_check\_regs} then, per operand, classifies the raw bits
using the IEEE-754 bit-pattern predicates in \texttt{fp\_utils.cuh}
(\texttt{half\_is\_nan}, \texttt{float\_is\_inf}, \texttt{double\_is\_subnorm},
etc. --- pure exponent/mantissa bit masking, no floating-point arithmetic
at all, so this cannot itself introduce rounding artifacts):

\begin{lstlisting}[language=C++]
uint32_t float_is_nan(uint32_t f) {
  return (float_exp(f) == 0xFF && float_mant(f) != 0) ? E_NAN : 0;
}
uint32_t float_is_inf(uint32_t f) {
  return (float_exp(f) == 0xFF && float_mant(f) == 0) ? sgninf(f) : 0;
}
\end{lstlisting}

The five exception bits are \texttt{E\_NAN}, \texttt{E\_INF},
\texttt{E\_NINF} (negative infinity, reported separately from positive),
\texttt{E\_SUB} (subnormal), and \texttt{E\_DIV0} --- the last of which is
really just ``this operand is exactly zero'', only checked for operands
the instruction table flags with \texttt{"div0": true} (i.e.\ divisor
operands of division-shaped instructions). \texttt{half2}/\texttt{bf162}
variants unpack two packed half-precision lanes from one 32-bit register
before classifying each.

The per-lane exception bitmask is then combined across the whole warp with
a \texttt{\_\_shfl\_sync} reduction (any lane's exception counts for the
warp), and only the warp's first active lane does the recording, to avoid
32-way redundant writes:

\begin{lstlisting}[language=C++]
if (first_laneid == laneid) {
  for (int op = 0; op < OPERANDS; op++) {
    uint32_t exce = exces[op];
    if (exce != 0) {
      uint32_t num_exceptions = recorder.record(inst_id, exce, op);
      if (num_exceptions == 0) {
        // first time this (instruction, operand) pair has ever
        // hit this exact exception -- push a live report
        for (int skip = 0; skip < 2; skip++) {
          exception_info ei(get_ctaid(), get_warpid(), inst_id, exce, op, UNKNOWN, skip);
          channel_dev->push(&ei, sizeof(exception_info));
        }
      }
    }
  }
}
\end{lstlisting}

\texttt{recorder.record} (\texttt{device\_recorder::record},
\texttt{recording.h}) is an \texttt{atomicAdd} into a flat device array,
indexed by \texttt{(instruction id, operand, exception bitmask)}. This is
the counter that \textbf{every} occurrence increments, whether or not a
message is pushed anywhere --- it is the source of the ``repeats'' numbers
in the final report. The channel push, by contrast, only happens the
\emph{first} time a given \texttt{(instruction, operand, exception-bits)}
triple is ever seen (\texttt{num\_exceptions == 0} means this was the very
first increment) --- this is what produces the live
\texttt{\#nixnan: error [\ldots] detected in operand \ldots} lines this
session saw streamed to the console, and it happens exactly once per
distinct site, no matter how many times that site actually fires.

\subsubsection*{The channel and receiver thread}

NVBit tools cannot easily do host-side I/O from device code, so
\texttt{nixnan.so} uses NVBit's lock-free \texttt{ChannelDev}/\texttt{ChannelHost}
ring buffer: device code \texttt{push}es fixed-size \texttt{exception\_info}
records; a dedicated host \texttt{std::thread} (\texttt{recv\_thread\_fun}
in \texttt{nixnan.cu}) polls \texttt{channel\_host.recv} and, for each
record, looks up the recorded instruction's SASS text, function name, and
source location (all previously stashed by \texttt{recorder::mk\_entry} at
instrumentation time) and prints the formatted error line. Every push
happens \emph{twice}, with a ``skip'' flag on the second copy --- this is
an NVBit idiom to guarantee the channel has fully drained even if the
consuming thread is mid-poll (\texttt{to\_skip()} entries are silently
discarded by the receiver).

\subsubsection*{The final aggregate report}

At context teardown (\texttt{nvbit\_at\_ctx\_term}), the device-side error
array is copied back to the host and folded into per-\texttt{(format,
is\_memory)} totals. For each of the five exception kinds, the code
computes \emph{two} numbers from the same per-instruction/operand counter
table:

\begin{lstlisting}[language=C++]
if (exce & error) {
  std::get<0>(exception_counts[{type, is_mem}][index]) += errors;   // total occurrences (repeats, so far)
  if (errors > 0) std::get<1>(exception_counts[{type, is_mem}][index])++;  // one more distinct site
}
...
// later: subtract the per-format unique-site count out of the running total
std::get<0>(exception_counts[{ftype, is_mem}][i]) -= std::get<1>(exception_counts[{ftype, is_mem}][i]);
\end{lstlisting}

That subtraction is exactly why the printed report reads
\texttt{NaN: 22 (154 repeats)}: 22 is the number of distinct
(instruction, operand) sites that ever produced a NaN even once, and 154
is \emph{additional} occurrences beyond each site's first (i.e.\ total
occurrences minus the 22 ``first hits'' already counted once each in the
unique-site tally) --- so \texttt{unique + repeats} is the true total
number of times that exception fired, summed over every site. This is
the exact metric this session used to compare \texttt{rd\_nixnan.cu}
against its 10x-grid-point variant: unique-site count is a static
property of the kernel's source code and stays fixed as the problem
scales, while repeats measures how much bad floating-point behavior an
actual run produced, and that number does scale with problem size.

\subsection{Memory-instruction checking}
\label{sec:memcheck}

If \texttt{INSTR\_MEM=1}, every global-memory store (\texttt{STG}, any
width) additionally gets a call to \texttt{nixnan\_check\_nans}
(\texttt{memcheck.cu}) inserted \texttt{IPOINT\_BEFORE}.
\texttt{meminstrumentation.cu} first has to work out \emph{what type} is
being stored, since a raw \texttt{STG} instruction's operand is just an
untyped register: it walks backwards through the current basic block
looking for the most recent instruction that wrote that register and
reuses \emph{that} instruction's already-known operand type
(\texttt{find\_type}). If no typed producer is found, the store is
skipped rather than guessed at. Unlike the exception path above, the
memory-instruction checker only ever reports \texttt{NaN} (not
Inf/Subnormal/Div0) --- this is why every ``\texttt{\ldots Memory
Operations}'' block in the report only ever prints a \texttt{NaN:} line.

\subsection{The exponent-histogram / bin-threshold mechanism}
\label{sec:histogram}

This is the machinery behind \texttt{HISTOGRAM}/\texttt{BIN\_SPEC\_FILE},
and it is what \texttt{autonixnan}'s ``extreme-value scan'' phase
(Section~\ref{sec:phase2}) is actually built on. It runs as a
\emph{parallel, independent} instrumentation pass alongside the exception
checker above --- every FP instruction that gets a
\texttt{nixnan\_check\_regs} call also, unconditionally, gets a call to
\texttt{nixnan\_fp\_histogram\_counter} inserted (\texttt{fp-histogram.cu});
whether that call does anything is gated at run time by
\texttt{histogram\_enabled} (true if \texttt{HISTOGRAM=1} or
\texttt{BIN\_SPEC\_FILE} is set).

\texttt{BIN\_SPEC\_FILE} is a small JSON document naming, per format, a
list of \texttt{[lower, upper]} \emph{unbiased signed exponent} ranges to
watch, plus a \texttt{count} (report every \texttt{count}-th value landing
in a watched bin) and optional \texttt{max\_reports} (self-terminate a
kernel-function's reporting after this many). If the named file does not
exist, \texttt{nixnan.so} writes an empty template there and
\texttt{exit(0)}s --- by design, the first invocation with a
not-yet-existing spec file is a no-op setup step, not a real run.

At run time, \texttt{record()} (in \texttt{fp-fun.cu}) does two things
for every FP value seen: it always \texttt{atomicAdd}s into a flat
\texttt{[format][exponent]} histogram array (this is what powers the
``Exponent range for f16: [-14, 15]'' summary line printed at
\texttt{term()}, when \texttt{HISTOGRAM=1} is used directly rather than
through \texttt{autonixnan}), and, if that exponent falls in one of the
function's configured bins, increments a per-bin counter and --- every
\texttt{count}-th time --- pushes a \texttt{bin has reached threshold}
report through its own separate channel/receiver-thread pair
(structurally identical to the exception path, but entirely separate
state). \textbf{Crucially, the valid exponent range for each format
explicitly excludes zero and infinity} (see \texttt{FP\_FORMAT\_RANGES} in
\texttt{bin/autonixnan}, and \texttt{bin\_from\_json}'s range validation in
\texttt{fp-histogram.cu}, which rejects any bin touching the
all-zeros/all-ones exponent encodings) --- \textbf{a value that jumps
straight from a normal, finite magnitude to \texttt{Inf} in a single
instruction never lands in any bin at all, and is therefore invisible to
this mechanism.} This is exactly what this session found empirically
with the mixed-precision-Cholesky Phase-D overflow experiment
(Section~\ref{sec:limitations}): the scan flagged an unrelated fp64 QR
kernel while missing the real, program-confirmed FP16 overflow entirely,
because that overflow was a one-step jump to \texttt{Inf} with no
gradually-climbing finite exponent trail leading up to it.

\subsection{\texttt{LOG\_KERNELS} mode}
\label{sec:logkernels}

As established above, \texttt{LOG\_KERNELS} disables all instrumentation.
What it \emph{does} do: at kernel-exit time,
\texttt{nvbit\_at\_cuda\_event} prints
\texttt{Kernel [name] execution time: N microseconds} for every launch
(gated on \texttt{time\_kernels || kernel\_logging\_enabled}, so this
line's presence does not actually require \texttt{TIME\_KERNELS=1} once
\texttt{LOG\_KERNELS} is set --- \texttt{autonixnan}'s
\texttt{collect\_kernels()} relies on exactly this). Separately,
\texttt{instrument\_function} also prints one
\texttt{Function [related] of kernel [top-level]} line for \emph{every}
NVBit-enumerated related function of a kernel, which for a library-heavy
binary (cuBLAS/cuSOLVER template instantiations) can be enormous --- an
n=4096 cuSOLVER Cholesky factorization run in this session produced over
10{,}000 such lines, none of which match \texttt{autonixnan}'s
\texttt{\#nixnan: Kernel [\ldots] execution time:} parser regex and are
therefore silently ignored by \texttt{\_parse\_kernel\_log}. They are
harmless noise from \texttt{autonixnan}'s point of view, but they do
inflate the log file and are worth knowing about if reading a
\texttt{LOG\_KERNELS} log by hand.

\subsection{The (currently) dead \texttt{NIXNAN\_TIMEOUT}}
\label{sec:timeout-dead}

\texttt{src/timer.h}/\texttt{timer.cu} implement a self-contained,
graceful internal timeout: \texttt{nixnan::timer::init()} reads
\texttt{NIXNAN\_TIMEOUT} (default 0, meaning disabled), and
\texttt{nixnan::timer::tool\_init(ctx)}, if the timeout is positive,
spawns a background thread that sleeps for that many seconds and then
calls \texttt{nnterminate("Timeout reached.")} --- a graceful,
in-process \texttt{SIGINT} to itself, flushing the log first.

\textbf{This code is never called.} A full read of \texttt{nixnan.cu}'s
\texttt{nvbit\_at\_init}, \texttt{nvbit\_tool\_init}, and
\texttt{nvbit\_at\_ctx\_term} --- the only three places that call into
any other \texttt{namespace}'s \texttt{init}/\texttt{tool\_init}/\texttt{term}
--- shows calls into \texttt{nixnan::fp\_histogram::\{init,tool\_init,term\}}
only. There is no call anywhere in the tree into
\texttt{nixnan::timer::\{init,tool\_init,term\}}. This was confirmed two
independent ways during this session:

\begin{itemize}
\item \texttt{grep -rn "timer::" src/} outside of \texttt{timer.cu} itself
  returns nothing.
\item The \texttt{GET\_VAR\_INT} macro used to read \texttt{NIXNAN\_TIMEOUT}
  \emph{always} prints a banner line when it runs, regardless of the
  value found. Every banner captured this session (which comes from
  \texttt{nvbit\_at\_init}, run unconditionally) lists
  \texttt{INSTR\_BEGIN} through \texttt{FUNCTION\_BLACKLIST} --- and never
  \texttt{NIXNAN\_TIMEOUT}. If the variable were actually being read, its
  banner line would be there.
\end{itemize}

The practical consequence: \texttt{bin/autonixnan}'s
\texttt{find\_extreme\_values()} sets \texttt{NIXNAN\_TIMEOUT} and
documents (in its own docstring) that this is what lets ``the target's
own timer, rather than the hard kill here'' end a long-running scan
gracefully. In the currently built \texttt{nixnan.so}, that never
happens --- \texttt{NIXNAN\_TIMEOUT} is silently ignored, and \emph{every}
timeout this session observed (``\texttt{Error: target program exceeded
Ns timeout}'') was Python's own \texttt{subprocess.run(timeout=\ldots)}
hard-killing the child with \texttt{SIGKILL} after the full
\texttt{args.timeout + TIMEOUT\_GRACE} (30s) window, not a graceful
internal exit. This does not make \texttt{autonixnan} incorrect --- the
external kill still bounds run time correctly --- but it does mean the
30-second grace period is not, in practice, ``slack given in case the
internal timer is a little late''; it is dead weight, since nothing
internal is ever going to fire first.

\section{Part II --- \texttt{bin/autonixnan}}

\texttt{autonixnan} runs three phases against the same target command,
in sequence, unconditionally launching the target process fresh each
time (there is no persistent daemon or warm cache between phases).

\subsection{Phase 1: \texttt{collect\_kernels()} --- the baseline}

\begin{lstlisting}[language=Python]
env["LD_PRELOAD"] = lib_path
env["LOG_KERNELS"] = log_file
elapsed = run_with_env(program_args, timeout, env)
\end{lstlisting}

This runs the target with instrumentation fully disabled (Section
\ref{sec:logkernels}), so wall time here is close to native (modulo the
overhead of NVBit's own related-function enumeration and the
\texttt{"Function [\ldots] of kernel [\ldots]"} print spam for
library-heavy binaries, which this session measured at 10-17s for a
cuSOLVER-based binary that runs in under 2s uninstrumented). The
\texttt{-t}/\texttt{--timeout} value is used \emph{directly, with no
grace period}, as the hard \texttt{subprocess.run} timeout for this
phase --- a target that is merely slow to start (library init, matrix
setup) can therefore fail here even though nothing is actually being
instrumented, which this session hit directly with a cuSOLVER-heavy
binary whose factorization setup alone exceeded a naively small
\texttt{-t}.

The resulting log is parsed by \texttt{\_parse\_kernel\_log}, which only
matches lines starting \texttt{"\#nixnan: Kernel ["} --- ignoring the
\texttt{"Function [\ldots] of kernel [\ldots]"} noise described above ---
extracting the kernel name and, if present, its microsecond execution
time, and building both an in-order call sequence and a
name~$\to$~\{count, total\_time\} table.

\subsection{Report assembly (between phases 1 and 2)}

Before running Phase 2, \texttt{main()} prints the ``Unique Kernels''
table (sorted by total time descending, truncated to the top 25 when
there are more), the baseline wall-clock time, and the full kernel call
sequence --- elided to the first and last 25 calls with a
\texttt{"(\ldots calls omitted)"} marker in between when there are more
than 50 total, so an n=4096 factorization's 900+ launches do not flood
the terminal.

\subsection{Phase 2: \texttt{find\_extreme\_values()} --- the exponent scan}
\label{sec:phase2}

\begin{lstlisting}[language=Python]
bs = _make_binspec(1, max_reports)   # count=1: report every single hit, capped at max_reports/kernel
env["BIN_SPEC_FILE"] = spec_file
env["LOGFILE"] = log_file
env["NIXNAN_TIMEOUT"] = str(timeout)      # see Sec 2.9 -- currently a no-op
env["INSTRUMENT_EXCEPTIONS"] = "0"        # exception checker off; histogram checker still runs
elapsed = run_with_env(program_args, timeout + TIMEOUT_GRACE, env, check=False)
\end{lstlisting}

\texttt{\_make\_binspec} auto-generates a \texttt{BIN\_SPEC\_FILE} that
watches the top 5\% and bottom 5\% of \emph{every} format's valid
(finite, non-zero) exponent range, for \emph{every} function --- since
the spec file has no per-function key, \texttt{fp-histogram.cu}'s
\texttt{process\_bin\_spec} applies the same global spec to whichever
function it is asked about, so in effect every instrumented kernel is
watched. This is a real, full NVBit exception-instrumentation-grade pass
over the entire binary (with the exception checker itself switched off
via \texttt{INSTRUMENT\_EXCEPTIONS=0}, leaving only the histogram
counter active) and is therefore subject to the same heavy per-distinct-kernel
overhead as any other full instrumentation run --- this session measured
it taking well over a minute for a cuSOLVER/cutlass-heavy binary with
hundreds of distinct kernel templates, versus milliseconds for a
handful of custom kernels called thousands of times.

\texttt{run\_with\_env} is called with \texttt{check=False}, so a
non-zero exit (which is the \emph{expected}, normal outcome here --- the
target hits \texttt{max\_reports} for some kernel and \texttt{nnterminate}s
itself with \texttt{SIGINT}) is not treated as an error; only an actual
\texttt{subprocess.TimeoutExpired} (after \texttt{timeout + 30}s) is
fatal to this phase.

The resulting log is scanned by \texttt{BIN\_REACHED\_RE} for
\texttt{"\#nixnan: <fmt> bin has reached threshold: function=\ldots
range=[lo,hi] count=N"} lines, tallied per function name, and the single
function with the most such reports across the whole run is chosen as
``the'' function of interest.

\subsection{Phase 3: \texttt{run\_instrumented()} --- the deep dive}

If Phase 2 found anything, \texttt{autonixnan} runs the target one final
time, with \texttt{LD\_PRELOAD} set and \emph{no other environment
variable} beyond a \texttt{FUNCTION\_WHITELIST} file containing exactly
the one function name Phase 2 flagged:

\begin{lstlisting}[language=Python]
env["LD_PRELOAD"] = lib_path
if whitelist_file:
    env["FUNCTION_WHITELIST"] = whitelist_file
elapsed = run_with_env(program_args, timeout, env, quiet=False)
\end{lstlisting}

Because no other variables are set, this run gets \texttt{INSTRUMENT\_EXCEPTIONS}'s
\emph{default} of 1 (full exception checking, Section~\ref{sec:exceptions})
and \texttt{HISTOGRAM}'s default of 0 (no histogram) --- but scoped, via
the whitelist, to \emph{only} the one function Phase 2 flagged (plus, per
\texttt{should\_instrument}, any function whose \emph{exact} name appears
in that whitelist; unrelated kernels are left completely uninstrumented,
running at native speed). \texttt{quiet=False} means this is the one
phase whose child process's \texttt{stdout}/\texttt{stderr} are not
suppressed --- this is why the live \texttt{\#nixnan: error [\ldots]}
lines and the final \texttt{------------ nixnan Report -----------}
aggregate (Section~\ref{sec:exceptions}) appear directly in
\texttt{autonixnan}'s own output during this phase and nowhere else.

This phase's own log file is opened and unconditionally
\texttt{os.unlink}ed in a \texttt{finally} block \emph{without ever being
read} --- \texttt{LOGFILE} is never set for this call (it is commented out
in the source), so the target's \texttt{nnout()} output goes to
\texttt{stderr}, which \texttt{quiet=False} already passes through to the
console directly; the temp file exists only so \texttt{tempfile.mkstemp}
has something to clean up, and \texttt{run\_instrumented}'s returned
\texttt{kernels\_seq}/\texttt{kernels\_unique} are consequently always
empty and are not used by \texttt{main()}.

\textbf{This is also the phase where the exponent-scan/exception-checker
scope mismatch documented in Section~\ref{sec:histogram} bites hardest}:
if the real fault is a hard jump straight to \texttt{Inf} (rather than a
gradual climb through extreme-but-finite exponents), Phase 2 will not
flag the kernel where it actually happens, so Phase 3's whitelist will
scope full exception-instrumentation to the \emph{wrong} kernel, and the
final aggregate report can show all-zero counts in every format even
though the target program's own internal exception detection (if it has
any) clearly fired. Section~\ref{sec:limitations} gives the exact
reproduction from this session.

\subsection{Timeout semantics, summarized}

\begin{center}
\begin{tabular}{lll}
\toprule
Phase & Hard kill after & Notes \\
\midrule
1: baseline & \texttt{-t} exactly & No grace period; catches slow startup too. \\
2: exponent scan & \texttt{-t + 30s} & Grace period, currently pointless (Sec.~\ref{sec:timeout-dead}). \\
3: deep dive & \texttt{-t} exactly & Same value re-used; not separately configurable. \\
\bottomrule
\end{tabular}
\end{center}

Every phase re-uses the \emph{same} \texttt{-t} value from the command
line; there is no way to give the baseline phase a short budget and the
scan phase a long one from the CLI as it stands.

\section{Empirically Observed Limitations}
\label{sec:limitations}

These are not code-reading conclusions but things directly reproduced
against real binaries this session, cross-checked against the source
above.

\begin{enumerate}
\item \textbf{Hard overflow jumps are invisible to the extreme-value
  scan.} Running a mixed-precision Cholesky iterative-refinement solver
  (\texttt{SC26-expts/mpir\_gpu}) with its documented overflow-forcing
  flags (\texttt{--noscale --bscale 1e5}) produced a program-confirmed
  FP16 overflow (\texttt{nonfinite=256} from the target's own internal
  check) on the very first iteration. \texttt{autonixnan}'s Phase 2 never
  flagged the kernel responsible; it flagged an unrelated fp64 QR kernel
  used only during test-matrix generation, and Phase 3's whitelisted
  full-instrumentation pass consequently reported zero FP16 exceptions.
  Running \texttt{nixnan.so} directly (no \texttt{autonixnan}, no
  whitelist, default \texttt{INSTRUMENT\_EXCEPTIONS=1}) against the same
  invocation found the real exceptions immediately --- filed under the
  \texttt{FP32} bucket, not \texttt{FP16}, because that solver
  configuration only rounds values into fp16's representable range via a
  cast; the arithmetic that actually overflows executes as native FP32
  SASS instructions. Neither of these is a nixnan bug: the histogram
  mechanism's zero/inf exclusion is by design (Section~\ref{sec:histogram}),
  and the FP32 attribution is simply an accurate report of which SASS
  instruction type executed on the bad value.
\item \textbf{The dead internal timeout} (Section~\ref{sec:timeout-dead})
  means every \texttt{-t} this session used was enforced purely
  externally, with a pointless 30s grace period tacked onto the scan
  phase.
\item \textbf{\texttt{LOG\_KERNELS} noise scales with library surface
  area, not with actual kernel diversity.} A cuSOLVER/cuBLAS-heavy binary
  produced 10{,}000+ log lines from related-function enumeration alone,
  none of which affect \texttt{autonixnan}'s parsed output but which do
  make the raw log file unwieldy to read by hand.
\end{enumerate}

\section{Practical Summary}

To reproduce any given phase by hand outside \texttt{autonixnan}, the
minimal commands are:

\begin{lstlisting}[language=bash]
# Phase 1 equivalent -- kernel catalog, no instrumentation
LD_PRELOAD=./nixnan.so LOG_KERNELS=/tmp/k.log ./target [args...]

# Phase 2 equivalent -- exponent-extreme scan, exceptions off
LD_PRELOAD=./nixnan.so BIN_SPEC_FILE=/tmp/spec.json LOGFILE=/tmp/scan.log \
  INSTRUMENT_EXCEPTIONS=0 ./target [args...]

# Phase 3 equivalent, or a direct full exception run (no whitelist
# needed if the binary is small/fast enough to instrument in full)
LD_PRELOAD=./nixnan.so ./target [args...]
\end{lstlisting}

The last form -- direct, unwhitelisted \texttt{nixnan.so}, no
\texttt{autonixnan} involved -- is the one to reach for whenever a target
is small enough to instrument in full and the exponent-scan's blind spot
(hard overflow jumps) is a concern, since it is the only one of the three
that does not depend on Phase 2 having correctly identified the kernel of
interest.

\end{document}
